\section{Einleitung}
\label{sec:einleitung}

\subsection{Projektziel}
SignalReport ist eine Open-Source-Anwendung zur kontinuierlichen Überwachung der Internet-Qualität. Im Gegensatz zu kommerziellen Tools ist es:

\begin{itemize}
    \item \textbf{Kostenlos} und Open Source
    \item \textbf{Einfach zu installieren} (kein Server nötig)
    \item \textbf{Plattformunabhängig} (Windows, macOS, Linux, NAS, Docker)
    \item \textbf{Datensicherheit} (alle Daten lokal gespeichert)
\end{itemize}

\subsection{Problemstellung}
Internetprobleme werden oft erst wahrgenommen, wenn sie bereits existieren. SignalReport ermöglicht:

\begin{itemize}
    \item Frühzeitige Erkennung von Qualitätsproblemen
    \item Historische Auswertung (z. B. "Warum ist das Netz immer um 19 Uhr langsam?")
    \item Nachweis für Internetanbieter bei Problemen
    \item Langfristige Dokumentation für Provider-Beschwerden
\end{itemize}

\subsection{Projektgeschichte}
SignalReport entstand zwischen Jänner und April 2026 als Abschlussprojekt des Diplomlehrgangs \enquote{Software Developer:in} (Kurs 18195015) am WIFI Wien. Die ursprüngliche Version wurde im April 2026 erfolgreich präsentiert und bewertet. Seither wird das Projekt als persönliches Open-Source-Vorhaben aktiv weiterentwickelt -- mit zusätzlichen Funktionen, robusteren Architekturentscheidungen und kontinuierlicher Pflege.

Diese Dokumentation beschreibt den \textbf{aktuellen Entwicklungsstand}, nicht den exakten Stand zum Zeitpunkt der Diplomprüfung. Wer den ursprünglichen, geprüften Stand begutachten möchte, findet ihn als Release \texttt{V1} im GitHub-Repository unter \newline\href{https://github.com/matthili/SignalReport/releases}{https://github.com/matthili/SignalReport/releases}.

\subsection{Projektumfang}
Der aktuelle Stand umfasst:
\begin{itemize}
    \item Java-Anwendung mit Javalin (Web-UI) und embedded H2-Twin-Datenbank
    \item Mess-Engines für Ping, DNS und HTTP (plattformspezifisches Ping)
    \item Statistische Auswertung (95. Perzentil, Jitter, Paketverlust)
    \item Live-Visualisierung mit Chart.js und Stunden-Heatmap
    \item PDF-Export (24h / 7\,Tage / 12\,Monate) und CSV-Export
    \item IP-Tracking für dynamische IP-Adressen
    \item Konfigurierbare Messziele und Maintenance-Fenster
    \item Challenge-Response-Authentifizierung (SHA-256, ohne HTTPS-Pflicht)
    \item Web-basierter Setup-Wizard (keine CLI nötig)
    \item Crash-resistente Datenpersistenz durch Twin-Datenbank-Architektur (eingeführt nach Diplomprüfung)
\end{itemize}